\documentclass[a4paper,12pt]{article}
\usepackage[utf8]{inputenc}
\usepackage{amsmath}
\usepackage{graphicx}
\usepackage{float}
\usepackage{hyperref}
\usepackage{amsfonts}
\usepackage{amssymb}
\usepackage{geometry}
\usepackage{listings}
\usepackage{xcolor}

\geometry{a4paper, margin=1in}

\title{Exercícios Práticos de Python: Decomposição SVD e Compressão de Imagens}
\author{Manus AI}
\date{Outubro de 2025}

\begin{document}

\maketitle

\section*{Introdução}
Esta lista de exercícios foi desenvolvida para complementar a aula sobre decomposição SVD e o teorema da melhor aproximação de posto $r$. Os exercícios progridem em dificuldade e abordam diferentes aspectos da aplicação prática do SVD em processamento de imagens.

\section*{Exercício 1: Implementação Básica do SVD (Nível: Iniciante)}

\textbf{Objetivo:} Familiarizar-se com a função \texttt{np.linalg.svd()} e verificar as propriedades básicas da decomposição.

\textbf{Enunciado:}
Crie uma função \texttt{analisar\_svd(matriz)} que:
\begin{enumerate}
    \item Receba uma matriz $A$ como entrada.
    \item Calcule a decomposição SVD: $A = U \times \Sigma \times V^T$.
    \item Imprima as dimensões de $U$, $\Sigma$ e $V^T$.
    \item Verifique se a reconstrução $A_{\text{reconstruída}} = U \cdot \text{diag}(s) \cdot V^T$ é igual à matriz original (use \texttt{np.allclose()}).
    \item Retorne o número de valores singulares não-nulos (posto da matriz).
\end{enumerate}

\textbf{Teste sua função com:}
\begin{itemize}
    \item Uma matriz $4 \times 3$ aleatória.
    \item Uma matriz identidade $5 \times 5$.
    \item Uma matriz de posto 2 (crie multiplicando duas matrizes de dimensões apropriadas).
\end{itemize}

\section*{Exercício 2: Carregamento e Análise de Imagem (Nível: Iniciante)}

\textbf{Objetivo:} Aprender a trabalhar com imagens como matrizes e analisar suas propriedades.

\textbf{Enunciado:}
Desenvolva um programa que:
\begin{enumerate}
    \item Carregue uma imagem em escala de cinza usando \texttt{matplotlib.pyplot.imread()} ou \texttt{PIL.Image.open()}.
    \item Converta para escala de cinza se necessário.
    \item Exiba informações sobre a imagem:
    \begin{itemize}
        \item Dimensões (altura $\times$ largura).
        \item Tipo de dados (\texttt{dtype}).
        \item Valores mínimo e máximo dos pixels.
        \item Posto da matriz da imagem.
    \end{itemize}
    \item Plote a imagem original usando \texttt{plt.imshow()}.
    \item Calcule e exiba o histograma dos valores dos pixels.
\end{enumerate}

\textbf{Dica:} Use \texttt{plt.imread()} para carregar e \texttt{plt.imshow(cmap='gray')} para visualizar em escala de cinza.

\section*{Exercício 3: Aproximação de Posto Baixo - Análise Quantitativa (Nível: Intermediário)}

\textbf{Objetivo:} Implementar o teorema da melhor aproximação de posto $r$ e analisar o erro de reconstrução.

\textbf{Enunciado:}
Crie uma função \texttt{aproximacao\_posto\_r(imagem, r)} que:
\begin{enumerate}
    \item Calcule a decomposição SVD da imagem.
    \item Construa a aproximação de posto $r$ usando apenas os primeiros $r$ valores singulares.
    \item Calcule o erro de Frobenius de duas formas:
    \begin{itemize}
        \item \textbf{Método direto:} $\|A - A_r\|_F$ (norma da diferença das matrizes).
        \item \textbf{Método teórico:} $\sqrt{\sigma_{r+1}^2 + \sigma_{r+2}^2 + \dots + \sigma_n^2}$ (usando valores singulares).
    \end{itemize}
    \item Verifique se os dois métodos produzem o mesmo resultado.
    \item Retorne a imagem aproximada e o erro calculado.
\end{enumerate}

\textbf{Teste com diferentes valores de $r$:} 1, 5, 10, 20, 50, 100 e compare os erros.

\section*{Exercício 4: Visualização Comparativa de Aproximações (Nível: Intermediário)}

\textbf{Objetivo:} Criar visualizações para comparar diferentes níveis de aproximação.

\textbf{Enunciado:}
Desenvolva um programa que:
\begin{enumerate}
    \item Carregue uma imagem de sua escolha.
    \item Calcule aproximações de posto $r$ para $r = [1, 5, 10, 25, 50, 100]$.
    \item Crie um subplot com $2 \times 3$ imagens mostrando:
    \begin{itemize}
        \item A imagem original no primeiro subplot.
        \item As 5 aproximações nos demais subplots.
    \end{itemize}
    \item Para cada aproximação, inclua no título:
    \begin{itemize}
        \item O valor de $r$.
        \item O erro de reconstrução.
        \item A porcentagem de compressão: $\left(1 - \frac{r(m+n+1)}{mn}\right) \times 100$.
    \end{itemize}
    \item Salve a figura como \texttt{comparacao\_aproximacoes.png}.
\end{enumerate}

\textbf{Desafio adicional:} Adicione uma curva mostrando como o erro decresce com o aumento de $r$.

\section*{Exercício 5: Análise de Eficiência de Compressão (Nível: Intermediário)}

\textbf{Objetivo:} Analisar o trade-off entre qualidade e compressão.

\textbf{Enunciado:}
Implemente uma função \texttt{analisar\_compressao(imagem)} que:
\begin{enumerate}
    \item Calcule aproximações para $r$ variando de 1 até o posto máximo da imagem.
    \item Para cada $r$, calcule:
    \begin{itemize}
        \item Taxa de compressão: $\text{razao\_compressao} = \frac{r(m+n+1)}{mn}$.
        \item Erro relativo: $\text{erro\_relativo} = \frac{\|A - A_r\|_F}{\|A\|_F}$.
        \item PSNR (Peak Signal-to-Noise Ratio): $\text{PSNR} = 20 \log_{10}\left(\frac{255}{\text{RMSE}}\right)$.
    \end{itemize}
    \item Plote três gráficos:
    \begin{itemize}
        \item Erro relativo vs. $r$.
        \item Taxa de compressão vs. $r$.
        \item PSNR vs. Taxa de compressão.
    \end{itemize}
    \item Identifique o "ponto ótimo" onde se obtém boa qualidade com alta compressão.
\end{enumerate}

\begin{enumerate}
    \item \textbf{Enunciado:}
    Crie uma função \texttt{analisar\_svd(matriz)} que:
    \begin{enumerate}
        \item Receba uma matriz $A$ como entrada.
        \item Calcule a decomposição SVD: $A = U \times \Sigma \times V^T$.
        \item Imprima as dimensões de $U$, $\Sigma$ e $V^T$.
        \item Verifique se a reconstrução $A_{\text{reconstruída}} = U \cdot \text{diag}(s) \cdot V^T$ é igual à matriz original (use \texttt{np.allclose()}).
        \item Retorne o número de valores singulares não-nulos (posto da matriz).
    \end{enumerate}
    \textbf{Teste sua função com:}
    \begin{itemize}
        \item Uma matriz $4 \times 3$ aleatória.
        \item Uma matriz identidade $5 \times 5$.
        \item Uma matriz de posto 2 (crie multiplicando duas matrizes de dimensões apropriadas).
    \end{itemize}

    \item \textbf{Enunciado:}
    Desenvolva um programa que:
    \begin{enumerate}
        \item Carregue uma imagem em escala de cinza usando \texttt{matplotlib.pyplot.imread()} ou \texttt{PIL.Image.open()}.
        \item Converta para escala de cinza se necessário.
        \item Exiba informações sobre a imagem:
        \begin{itemize}
            \item Dimensões (altura $\times$ largura).
            \item Tipo de dados (\texttt{dtype}).
            \item Valores mínimo e máximo dos pixels.
            \item Posto da matriz da imagem.
        \end{itemize}
        \item Plote a imagem original usando \texttt{plt.imshow()}.
        \item Calcule e exiba o histograma dos valores dos pixels.
    \end{enumerate}

    \item \textbf{Enunciado:}
    Crie uma função \texttt{aproximacao\_posto\_r(imagem, r)} que:
    \begin{enumerate}
        \item Calcule a decomposição SVD da imagem.
        \item Construa a aproximação de posto $r$ usando apenas os primeiros $r$ valores singulares.
        \item Calcule o erro de Frobenius de duas formas:
        \begin{itemize}
            \item \textbf{Método direto:} $\|A - A_r\|_F$ (norma da diferença das matrizes).
            \item \textbf{Método teórico:} $\sqrt{\sigma_{r+1}^2 + \sigma_{r+2}^2 + \dots + \sigma_n^2}$ (usando valores singulares).
        \end{itemize}
        \item Verifique se os dois métodos produzem o mesmo resultado.
        \item Retorne a imagem aproximada e o erro calculado.
    \end{enumerate}
    \textbf{Teste com diferentes valores de $r$:} 1, 5, 10, 20, 50, 100 e compare os erros.

    \item \textbf{Enunciado:}
    Desenvolva um programa que:
    \begin{enumerate}
        \item Carregue uma imagem de sua escolha.
        \item Calcule aproximações de posto $r$ para $r = [1, 5, 10, 25, 50, 100]$.
        \item Crie um subplot com $2 \times 3$ imagens mostrando:
        \begin{itemize}
            \item A imagem original no primeiro subplot.
            \item As 5 aproximações nos demais subplots.
        \end{itemize}
        \item Para cada aproximação, inclua no título:
        \begin{itemize}
            \item O valor de $r$.
            \item O erro de reconstrução.
            \item A porcentagem de compressão: $\left(1 - \frac{r(m+n+1)}{mn}\right) \times 100$.
        \end{itemize}
        \item Salve a figura como \texttt{comparacao\_aproximacoes.png}.
    \end{enumerate}
    \textbf{Desafio adicional:} Adicione uma curva mostrando como o erro decresce com o aumento de $r$.

    \item \textbf{Enunciado:}
    Implemente uma função \texttt{analisar\_compressao(imagem)} que:
    \begin{enumerate}
        \item Calcule aproximações para $r$ variando de 1 até o posto máximo da imagem.
        \item Para cada $r$, calcule:
        \begin{itemize}
            \item Taxa de compressão: $\text{razao\_compressao} = \frac{r(m+n+1)}{mn}$.
            \item Erro relativo: $\text{erro\_relativo} = \frac{\|A - A_r\|_F}{\|A\|_F}$.
            \item PSNR (Peak Signal-to-Noise Ratio): $\text{PSNR} = 20 \log_{10}\left(\frac{255}{\text{RMSE}}\right)$.
        \end{itemize}
        \item Plote três gráficos:
        \begin{itemize}
            \item Erro relativo vs. $r$.
            \item Taxa de compressão vs. $r$.
            \item PSNR vs. Taxa de compressão.
        \end{itemize}
        \item Identifique o "ponto ótimo" onde se obtém boa qualidade com alta compressão.
    \end{enumerate}
    \textbf{Fórmula do RMSE:} $\sqrt{\text{mean}((A - A_r)^2)}$.
    \item \textbf{Enunciado:}
    Desenvolva um sistema de compressão para imagens coloridas:
    \begin{enumerate}
        \item Carregue uma imagem colorida (RGB).
        \item Separe os canais R, G e B.
        \item Para cada canal, aplique SVD e mantenha apenas os primeiros $r$ componentes.
        \item Reconstrua a imagem colorida combinando os canais comprimidos.
        \item Implemente uma função \texttt{comprimir\_imagem\_colorida(imagem, r\_red, r\_green, r\_blue)} que permita diferentes níveis de compressão para cada canal.
        \item Compare visualmente:
        \begin{itemize}
            \item Imagem original.
            \item Compressão uniforme (mesmo $r$ para todos os canais).
            \item Compressão adaptativa ($r$ diferente por canal, baseado na variância de cada canal).
        \end{itemize}
    \end{enumerate}
    \textbf{Dica:} O canal verde geralmente contém mais informação visual importante.

    \item \textbf{Enunciado:}
    Crie três métodos para determinar automaticamente o rank ótimo:
    \begin{enumerate}
        \item \textbf{Método 1 - Critério de Energia:}
        \begin{itemize}
            \item Mantenha componentes até que 99\% da "energia" (soma dos quadrados dos valores singulares) seja preservada.
        \end{itemize}
        \item \textbf{Método 2 - Critério do Cotovelo:}
        \begin{itemize}
            \item Encontre o ponto de inflexão na curva dos valores singulares usando a segunda derivada.
        \end{itemize}
        \item \textbf{Método 3 - Critério de Qualidade Visual:}
        \begin{itemize}
            \item Pare quando o PSNR atingir um limiar (ex: 30 dB).
        \end{itemize}
    \end{enumerate}
    Implemente uma função \texttt{rank\_otimo(imagem, metodo='energia', limiar=0.99)} e compare os resultados dos três métodos em diferentes tipos de imagem (fotos, desenhos, texturas).

    \item \textbf{Enunciado:}
    Conduza um estudo comparativo:
    \begin{enumerate}
        \item Colete 4 tipos diferentes de imagem:
        \begin{itemize}
            \item Foto com muitos detalhes (paisagem).
            \item Imagem com poucas cores (logo, desenho).
            \item Textura repetitiva (padrão).
            \item Imagem com gradientes suaves (céu, pôr do sol).
        \end{itemize}
        \item Para cada imagem, analise:
        \begin{itemize}
            \item Distribuição dos valores singulares (plot dos $\sigma_i$).
            \item Taxa de decaimento dos valores singulares.
            \item Rank efetivo (número de valores singulares > 1\% do máximo).
            \item Qualidade de reconstrução para diferentes valores de $r$.
        \end{itemize}
        \item Crie um relatório automatizado que:
        \begin{itemize}
            \item Gere gráficos comparativos.
            \item Identifique qual tipo de imagem comprime melhor.
            \item Explique os resultados baseado nas características visuais.
        \end{itemize}
    \end{enumerate}

    \item \textbf{Enunciado:}
    Desenvolva um programa interativo usando \texttt{matplotlib.widgets} que:
    \begin{enumerate}
        \item Permita carregar diferentes imagens.
        \item Tenha um slider para ajustar o valor de $r$ em tempo real.
        \item Mostre simultaneamente:
        \begin{itemize}
            \item Imagem original e comprimida lado a lado.
            \item Valores atuais de: $r$, erro, taxa de compressão, PSNR.
            \item Gráfico dos valores singulares com linha indicando o corte atual.
        \end{itemize}
        \item Inclua botões para:
        \begin{itemize}
            \item Salvar a imagem comprimida atual.
            \item Resetar para a imagem original.
            \item Alternar entre diferentes métricas de qualidade.
        \end{itemize}
    \end{enumerate}
    \textbf{Dica:} Use \texttt{matplotlib.widgets.Slider} e \texttt{plt.subplots\_adjust()} para layout.

    \item \textbf{Enunciado:} 
    Implemente um sistema de benchmark que:
    \begin{enumerate}
        \item Teste a decomposição SVD em imagens de diferentes tamanhos (64$\times$64, 128$\times$128, 256$\times$256, 512$\times$512, 1024$\times$1024).
        \item Meça e compare:
        \begin{itemize}
            \item Tempo de decomposição SVD.
            \item Tempo de reconstrução para diferentes valores de $r$.
            \item Uso de memória para armazenar $U$, $\Sigma$, $V^T$ vs. imagem original.
            \item Tempo de I/O para salvar/carregar dados comprimidos.
        \end{itemize}
        \item Investigue o impacto do tipo de dados (\texttt{float32} vs \texttt{float64}).
        \item Compare com outros métodos de compressão (JPEG) em termos de:
        \begin{itemize}
            \item Taxa de compressão.
            \item Qualidade visual (PSNR, SSIM).
            \item Tempo de processamento.
        \end{itemize}
        \item Gere gráficos mostrando como a performance escala com o tamanho da imagem.
    \end{enumerate}
    
\end{enumerate}\sqrt{\text{mean}((A - A_r)^2)}$.

\end{document}